% Auto-generated by generate_tex_fragments.py -- do not hand-edit.
% \bibitem text copied verbatim from paper/sections/refs_all.tex
% for every reference key this document cites.
\bibitem{ref1} GBD 2023 IPV and SVAC Collaborators. Disease burden attributable to intimate partner violence against females and sexual violence against children in 204 countries and territories, 1990-2023: a systematic analysis for the Global Burden of Disease Study 2023. Lancet (London, England). 2026. PMID: 41386261.
\bibitem{ref2} Sardinha L. Global, regional, and national prevalence estimates of physical or sexual, or both, intimate partner violence against women in 2018. Lancet (London, England). 2022. PMID: 35182472.
\bibitem{ref3} Kuswanto H. Prevalence of and factors associated with female child marriage in Indonesia. PLoS ONE. 2024. PMID: 38968277.
\bibitem{ref4} Rumble L. An empirical exploration of female child marriage determinants in Indonesia. BMC Public Health. 2018. PMID: 29587705.
\bibitem{ref5} Islam MK. Attitudes to and experiences of intimate partner violence among Rohingya women who married before eighteen years of age. Global Health Action. 2021. PMID: 34323166.
\bibitem{ref6} Hawkins AJ. Does marriage and relationship education work? A meta-analytic study. Journal of Consulting and Clinical Psychology. 2008. PMID: 18837590.
\bibitem{ref9} Markman HJ. A randomized clinical trial of the effectiveness of premarital intervention: moderators of divorce outcomes. Journal of Family Psychology. 2013. PMID: 23421844.
\bibitem{ref10} Williamson HC. Risk moderates the outcome of relationship education: a randomized controlled trial. Journal of Consulting and Clinical Psychology. 2015. PMID: 25664643.
\bibitem{ref16} Upadhyay Ushma D. Development and validation of a reproductive autonomy scale. Studies in Family Planning. 2014. PMID: 24615573.
\bibitem{ref17} McCauley Heather L. Psychometric properties and refinement of the Reproductive Coercion Scale. Contraception. 2017. PMID: 27639927.
\bibitem{ref18} Upadhyay Ushma D. Development and Validation of the Sexual and Reproductive Empowerment Scale for Adolescents and Young Adults. The Journal of Adolescent Health. 2021. PMID: 32690468.
\bibitem{ref19} Dagher Myriam. Adaptation and psychometric assessment of a sexual and reproductive empowerment scale in Arabic among refugee and non-refugee adolescent girls. BMC Medical Research Methodology. 2024. PMID: 39266993.
\bibitem{ref20} Alomair N. Factors influencing sexual and reproductive health of Muslim women: a systematic review. Reproductive Health. 2020. PMID: 32138744.
\bibitem{ref24} Abramsky T. Findings from the SASA! Study: a cluster randomized controlled trial to assess the impact of a community mobilization intervention to prevent violence against women and reduce HIV risk in Kampala, Uganda. BMC medicine. 2014. PMID: 25248996.
\bibitem{ref27} Shaw B. Shifting Norms in Faith Communities to Reduce Intimate Partner Violence: Results from a Cluster Randomized Controlled Trial in Nigeria. Journal of interpersonal violence. 2023. PMID: 37329160.
\bibitem{ref33} Hawkins AJ. Exploring programmatic moderators of the effectiveness of marriage and relationship education programs: a meta-analytic study. Behavior Therapy. 2012. PMID: 22304880.
\bibitem{ref38} O'Doherty L. Screening women for intimate partner violence in healthcare settings. Cochrane Database of Systematic Reviews. 2015. PMID: 26200817.
\bibitem{ref40} US Preventive Services Task Force. Screening for Intimate Partner Violence and Caregiver Abuse of Older or Vulnerable Adults: US Preventive Services Task Force Recommendation Statement. JAMA. 2025. PMID: 40553450.
\bibitem{ref42} Betron. Effectiveness of a clinic-based counselling intervention on risk of experiencing intimate partner violence and reproductive coercion: a matched-pair cluster-controlled trial in Ebonyi and Sokoto, Nigeria. BMJ Global Health. 2025. PMID: 41151831.
\bibitem{ref45} Karakurt. Couples Therapy for Intimate Partner Violence: A Systematic Review and Meta-Analysis. Journal of marital and family therapy. 2016. PMID: 27377617.
\bibitem{ref46} McCollum. Couples treatment for interpersonal violence: a review of outcome research literature and current clinical practices. Violence and victims. 2008. PMID: 18624089.
\bibitem{ref47} Schneider. From scared to repaired: using an attachment-based perspective to understand situational couple violence. Journal of marital and family therapy. 2014. PMID: 25039387.
\bibitem{ref49} Robertson. Women and men's use of coercive control in intimate partner violence. Violence and victims. 2011. PMID: 21780535.
\bibitem{ref51} Evans. "Even 'Daily' is Not Enough": How Well Do We Measure Domestic Violence and Abuse?-A Think-Aloud Study of a Commonly Used Self-Report Scale. Violence and victims. 2016. PMID: 26645540.
\bibitem{ref92} Taft. Examining strength at home couples to prevent intimate partner violence on a military installation: A randomized controlled trial. Journal of consulting and clinical psychology. 2024. PMID: 38206858.
\bibitem{ref93} Ghuman SJ. Women's autonomy and child survival: a comparison of Muslims and non-Muslims in four Asian countries. Demography. 2003. PMID: 12962056.
\bibitem{ref94} Chen. Intimate partner violence: counseling, community resources, and legal issues for IPV victims and perpetrators. FP essentials. 2013. PMID: 24053261.
\bibitem{ref95} Decker. Implementing Trauma-Informed Partner Violence Assessment in Family Planning Clinics. Journal of women's health. 2017. PMID: 28375750.
\bibitem{ref96} MacMillan. Approaches to screening for intimate partner violence in health care settings: a randomized trial. JAMA. 2006. PMID: 16882959.
\bibitem{ref97} Ford-Gilboe. Development of a brief measure of intimate partner violence experiences: the Composite Abuse Scale (Revised)-Short Form (CASR-SF). BMJ open. 2016. PMID: 27927659.
\bibitem{ref98} Hegarty. Validation of the Coercive-Composite Abuse Scale (C-CAS) in an Australian sample of women experiencing intimate partner abuse. BMJ open. 2026. PMID: 42323273.
\bibitem{ref99} Wilson. Psychometric Properties of the Coercion in Intimate Partner Relationships Scale. Assessment. 2023. PMID: 34167331.
\bibitem{ref100} Sarac A. Cultural adaptation and pilot evaluation of the Couple Learning Program (CLP) among young Turkish couples: A mixed-method study. Acta Psychologica. 2026. PMID: 42442188.
\bibitem{ref101} Welland C. Culturally specific treatment for partner-abusive Latino men: a qualitative study to identify and implement program components. Violence and Victims. 2010. PMID: 21287968.
\bibitem{ref102} Sikoki B. A qualitative study on perceptions of adolescents' sexual and reproductive health education in Yogyakarta, Indonesia. International Journal of Adolescent Medicine and Health. 2024. PMID: 39395200.
\bibitem{ref103} Fauk NK. Cultural and religious determinants of HIV transmission: A qualitative study with people living with HIV in Belu and Yogyakarta, Indonesia. PLoS ONE. 2021. PMID: 34780506.
\bibitem{ref104} Grisurapong S. Establishing a one-stop crisis center for women suffering violence in Khonkaen hospital, Thailand. International Journal of Gynaecology and Obstetrics. 2002. PMID: 12429436.
\bibitem{ref105} Nagae M. [Domestic violence as a women's health issue--activities of health professionals in Thailand]. Japanese Journal of Public Health (Nihon Koshu Eisei Zasshi). 2004. PMID: 15162975.
\bibitem{ref106} Williamson HC. Does premarital education decrease or increase couples' later help-seeking? Journal of Family Psychology. 2014. PMID: 24294929.
\bibitem{ref107} Williamson HC. Barriers and facilitators of relationship help-seeking among low-income couples. Journal of Family Psychology. 2019. PMID: 30489129.
\bibitem{ref108} Oduenyi. Integrating gender-based violence first-line response into family planning and antenatal care services: an assessment of feasibility, appropriateness and acceptability to healthcare providers in Sokoto and Ebonyi States, Nigeria. BMC Health Services Research. 2025. PMID: 41267027.
\bibitem{ref109} Rhodes. The anatomy of a community health center system-level intervention for intimate partner violence. Journal of Urban Health. 2014. PMID: 23917943.
\bibitem{ref110} Miller Elizabeth. A family planning clinic-based intervention to address reproductive coercion: a cluster randomized controlled trial. Contraception. 2016. PMID: 26892333.
\bibitem{ref111} Zaher. Effect of domestic violence training: systematic review of randomized controlled trials. Canadian Family Physician. 2014. PMID: 25022633.
\bibitem{ref112} Jacubowski N. Marriage is not a safe place: heterosexual marriage and HIV-related vulnerability in Indonesia. Culture, Health \& Sexuality. 2008. PMID: 18038283.
\bibitem{ref113} Michau L. SASA! Together: An evolution of the SASA! approach to prevent violence against women. Evaluation and program planning. 2021. PMID: 33578229.
\bibitem{ref114} Stern E. Lessons learned from implementing Indashyikirwa in Rwanda - an adaptation of the SASA! approach to prevent and respond to intimate partner violence. Evaluation and program planning. 2018. PMID: 30125773.
\bibitem{ref115} Clyde TL. Revising Premarital Relationship Interventions for the Next Generation. Journal of Marital and Family Therapy. 2020. PMID: 30725473.
\bibitem{ref116} Silliman B. Improving practice in marriage preparation. Journal of Sex \& Marital Therapy. 1999. PMID: 10081741.
\bibitem{ref118} Zust BL. 10-Year Study of Christian Church Support for Domestic Violence Victims: 2005-2015. Journal of Interpersonal Violence. 2021. PMID: 33729071.
\bibitem{ref160} Schenck D et al. Ethical Analysis and Policy Recommendations Regarding Domino Liver Transplantation. Transplantation. 2018; PMID: 29708521.
\bibitem{ref164} Messing JT et al. The average predictive validity of intimate partner violence risk assessment instruments. Journal of Interpersonal Violence. 2013; PMID: 23262817.
\bibitem{ref165} Wang PL et al. Assessing the Danger: Validation of Taiwan Intimate Partner Violence Danger Assessment. Journal of Interpersonal Violence. 2015; PMID: 25315482.
\bibitem{ref168} Gerth J et al. [The Ontario Domestic Assault Risk Assessment (ODARA) - Validity und authorised German translation of an intimate partner violence screening tool]. Fortschritte der Neurologie-Psychiatrie. 2014; PMID: 25383928.
\bibitem{ref171} Hogan NR et al. Investigating racial disparities in violence risk assessment using the Spousal Assault Risk Assessment Guide-Version 3. Psychological Assessment. 2024; PMID: 38512165.
